\section{Rev B-Gateway --- Gateway Node}
\label{sec:rev-b-gateway}

\begin{diagram}
            CAN/CAN-FD
                 |
             transceiver
                 |
                MCU
                 |
               SPI
                 |
              LAN8651
                 |
            10BASE-T1S
\end{diagram}

\subsection{MCU selection}

\begin{decision}[title={STM32G4-class MCU (e.g., STM32G474)}]
Integrated FDCAN peripherals give CAN-FD without an external controller IC (just a transceiver); ample SPI/UART/ADC/timer resources cover the LAN8651 SPI link plus CAN/LIN transceiver control; mainstream toolchain support and an inexpensive Nucleo board let firmware bring-up run in parallel with layout instead of waiting for Rev~B silicon. This is a prototype-phase choice only -- production Gateway MCU selection is out of scope (Section~\ref{sec:exec-summary}).
\end{decision}

\subsection{What Rev B-Gateway integrates}

\begin{itemize}[itemsep=0.1em]
\item CAN-FD (via the MCU's integrated FDCAN peripheral plus an external transceiver).
\item Optional LIN transceiver.
\item Hardware TX enable, safe default = disabled (Section~\ref{sec:gateway-modes}).
\item Watchdog, independent of MCU firmware health.
\item Protected 12/24\,V input, same protection philosophy as the E2B boards (Section~\ref{sec:platform-strategy}).
\item Rail voltage and current monitoring.
\item Debug/programming access (SWD).
\item LAN8651 for T1S access.
\end{itemize}

\subsection{Rev B pass gate (both architectures)}

Rev~B must prove \emph{both} customer architectures, not either in isolation:

\begin{diagram}
GREENFIELD                              BROWNFIELD

    Clicker + LAN8651                       CAN devices
          |                                     |
       T1S                                    Gateway
      /   \                                     |
LAN8660  LAN8661                             LAN8651
   |        |                                   |
actuator  lighting                             T1S
                                                 |
                                              Clicker
\end{diagram}

\begin{decision}[title={Rev B is the load-bearing deliverable}]
If both systems do not work reliably on Rev~B hardware, do not rush to Rev~C. Rev~C exists to make an already-proven architecture presentable (Section~\ref{sec:rev-c}); it cannot substitute for validation that did not happen on Rev~B.
\end{decision}
